% =====================================================================
%  2bat-ud2-14.tex  —  SESSIÓ 14: Repàs global i síntesi (estacions)
% =====================================================================
\horatitol{Sessió 14 — Repàs global i síntesi (estacions)}%
{Repàs actiu per quatre estacions, una per pas del mètode, i mapa final que connecta tot el procés de decisió.}

\subsection{Idea clau}

Repassem \textbf{tota la unitat} de manera activa, rotant per \textbf{quatre
estacions}, una per cada pas del mètode. Al final construirem un \textbf{mapa de la
unitat} que connecta els quatre passos: \textbf{traduir $\rightarrow$ representar
$\rightarrow$ optimitzar $\rightarrow$ interpretar}.

\begin{idea}
\textbf{Les quatre estacions}
\begin{enumerate}[leftmargin=1.6em, itemsep=2pt]
  \item \textbf{Traduir:} passar un enunciat social a un sistema d'inequacions.
  \item \textbf{Representar:} dibuixar la regió factible d'un sistema i trobar-ne els
        vèrtexs.
  \item \textbf{Optimitzar:} avaluar una funció objectiu als vèrtexs i triar el màxim
        o el mínim.
  \item \textbf{Interpretar i casos especials:} llegir una solució òptima, o reconèixer
        solució múltiple / regió no acotada.
\end{enumerate}
\end{idea}

\subsection{Les estacions}

\begin{exemple}[Estació 1 — traduir]
«Una entitat fa $x$ tallers i $y$ xerrades; com a màxim $10$ activitats en total, i com
a màxim $6$ tallers; res negatiu.» $\rightarrow$ $\;x+y\le 10,\;x\le 6,\;x\ge 0,\;y\ge 0$.
\end{exemple}

\begin{exemple}[Estació 3 — optimitzar]
Sobre els vèrtexs $(0,0)$, $(6,0)$, $(4,3)$, $(0,5)$, per maximitzar $P=2x+3y$
avaluem: $P(6,0)=12$, $P(4,3)=17$, $P(0,5)=15$. Màxim $P=17$ a $(4,3)$.
\end{exemple}

\subsection{Exercicis resolts}

\begin{exresolt}[model de l'Estació 2: representar i trobar vèrtexs]
Dibuixa la regió factible de $\;x+y\le 6,\;x+3y\le 12,\;x\ge 0,\;y\ge 0\;$ i dóna'n
els vèrtexs.
\solucio
Vèrtexs: $(0,0)$; eix $x$ (mana $x\le 6$) $\rightarrow (6,0)$; creuament $x+y=6$ i
$x+3y=12$ ($2y=6\Rightarrow y=3$, $x=3$) $\rightarrow (3,3)$; eix $y$ (mana $y\le 4$)
$\rightarrow (0,4)$.

\begin{center}
\begin{tikzpicture}
\begin{axis}[udaxis, width=8cm, height=4.8cm,
    xlabel={\small $x$}, ylabel={\small $y$},
    xmin=0, xmax=8, ymin=0, ymax=6, xtick={0,2,4,6}, ytick={0,2,4}]
  \fill[udblue!18] (axis cs:0,0) -- (axis cs:6,0) -- (axis cs:3,3) -- (axis cs:0,4) -- cycle;
  \addplot[udblue, domain=0:6, samples=2]{6-x};
  \addplot[udnavy, domain=0:8, samples=2]{(12-x)/3};
  \addplot[udnavy, only marks, mark=*, mark size=1.4pt] coordinates {(0,0)(6,0)(3,3)(0,4)};
\end{axis}
\end{tikzpicture}
\end{center}
\resposta{Vèrtexs $(0,0)$, $(6,0)$, $(3,3)$ i $(0,4)$.}
\end{exresolt}

\begin{exresolt}[model de l'Estació 4: interpretar i casos especials]
Sobre la regió anterior, maximitza $P=2x+2y$. Què observes?
\solucio
$P(0,0)=0$, $P(6,0)=12$, $P(3,3)=12$, $P(0,4)=8$. El màxim $P=12$ s'assoleix a
$(6,0)$ \emph{i} a $(3,3)$: és una \textbf{solució múltiple} (tot el segment entre
ells és òptim). Socialment, vol dir que hi ha diverses maneres igual de bones de
repartir els recursos.
\resposta{Solució múltiple: $P=12$ a $(6,0)$ i $(3,3)$, i a tot el segment entremig.}
\end{exresolt}

\subsection{Exercicis per fer}

\noindent\textit{Una tasca per estació. Rota i resol-les totes.}

\begin{exercicis}
\begin{exlist}
  \item \textbf{(Estació 1)} Tradueix a un sistema d'inequacions: «una associació
        reparteix $x$ lots de menjar i $y$ d'higiene; com a màxim $20$ lots en total i
        com a màxim $12$ de menjar; res negatiu».
  \item \textbf{(Estació 2)} Troba els vèrtexs de la regió $\;2x+y\le 8,\;x\ge 0,\;
        y\ge 0$.
  \item \textbf{(Estació 2)} Calcula el vèrtex on es creuen $x+y=7$ i $2x+y=10$.
  \item \textbf{(Estació 3)} Sobre els vèrtexs $(0,0)$, $(4,0)$, $(2,4)$, $(0,5)$,
        maximitza $P=3x+2y$.
  \item \textbf{(Estació 4)} Una regió és no acotada cap amunt, amb vèrtexs de baix
        $(0,5)$, $(2,2)$, $(6,0)$. Minimitza $C=3x+4y$. Té sentit buscar-ne el màxim?
  \item \textbf{(Mapa de la unitat)} Completa el mapa del mètode amb una frase per a
        cada pas:
        \begin{center}
        \renewcommand{\arraystretch}{1.4}
        \begin{tabular}{p{3.6cm} p{8.8cm}}
        \toprule
        \textbf{Pas} & \textbf{Què fem} \\
        \midrule
        1. Traduir & \rule{8cm}{0.4pt} \\
        2. Representar & \rule{8cm}{0.4pt} \\
        3. Optimitzar & \rule{8cm}{0.4pt} \\
        4. Interpretar & \rule{8cm}{0.4pt} \\
        \bottomrule
        \end{tabular}
        \end{center}
\end{exlist}
\end{exercicis}

% ---------------------------------------------------------------------
\fullsolucions{Sessió 14}
\begin{solucions}
\begin{sollist}
  \item $x+y\le 20$, \;$x\le 12$, \;$x\ge 0$, \;$y\ge 0$.
  \item Vèrtexs $(0,0)$, $(4,0)$ i $(0,8)$. (Sobre $y=0$: $2x\le 8\Rightarrow x\le 4$;
        sobre $x=0$: $y\le 8$.)
  \item Restant: $(2x+y)-(x+y)=10-7\Rightarrow x=3$; $y=7-3=4$. Vèrtex $(3,4)$.
  \item $P(0,0)=0$, $P(4,0)=12$, $P(2,4)=6+8=14$, $P(0,5)=10$. Màxim $P=14$ a $(2,4)$.
  \item $C(0,5)=20$, $C(2,2)=6+8=14$, $C(6,0)=18$. Mínim $C=14$ a $(2,2)$. El màxim
        \textbf{no} té sentit: la regió és no acotada i $C$ creix sense límit.
  \item Resposta oberta. Idees: 1. traduir l'enunciat a inequacions (variables +
        restriccions); 2. dibuixar la regió factible i trobar els vèrtexs; 3. avaluar
        la funció objectiu als vèrtexs i triar màxim/mínim; 4. interpretar l'òptim en
        termes reals i reconèixer casos especials.
\end{sollist}
\end{solucions}
